% ===========================================================================
\section{环的定义}\label{sec:rings}

\subsection{环与含幺环}

\begin{definition}[环（Ring）]
环 $(R, +, \cdot)$ 是一个集合 $R$ 连同两个二元运算——\keyword{加法} $+$
与\keyword{乘法} $\cdot$（常省略记 $ab$），满足：
\begin{enumerate}[leftmargin=1.6em,itemsep=2pt]
  \item $(R, +)$ 是\keyword{阿贝尔群}：结合律、交换律、存在零元 $0$、
        每个 $a$ 有负元 $-a$；
  \item $(R, \cdot)$ 是\keyword{半群}：结合律 $(ab)c = a(bc)$；
  \item \keyword{分配律}（乘法对加法）：
        \[
          a(b+c) = ab + ac, \qquad (a+b)c = ac + bc.
        \]
\end{enumerate}
\end{definition}

\begin{definition}[含幺环]
若存在 $1\in R$ 使对任意 $a\in R$ 有 $1\cdot a = a\cdot 1 = a$，则称 $R$
为\keyword{含幺环}。
\end{definition}

\begin{remark}[术语说明]
环的定义在文献中不统一：一派把不一定含幺的结构叫环（ring），含幺的叫
含幺环；另一派把含幺的叫环，不一定含幺的叫 \keyword{rng}（把 ``ring''
中的 ``i''（identity）丢掉）。后面真正关心的对象（如域）几乎都是含幺
且交换的。
\end{remark}

\begin{example}[无幺元的环]
在任一阿贝尔群上定义乘法 $a\cdot b = 0$（任意两元素相乘恒为零），则除
“存在幺元”外其余公理全部满足。
\end{example}

\subsection{交换环、除环、域}

\begin{definition}[交换环]
乘法满足交换律 $ab = ba$ 的环称为\keyword{交换环}。
\end{definition}

\begin{definition}[除环（Division Ring）]
若含幺环 $R$ 中每个非零元素都有乘法逆，即对任意 $a\ne 0$ 存在 $a^{-1}$
使
\[
  a\cdot a^{-1} = a^{-1}\cdot a = 1,
\]
则称 $R$ 为\keyword{除环}（体）。
\end{definition}

\begin{definition}[域（Field）]
交换的除环称为\keyword{域}。即：含幺、交换、每个非零元可逆。
\end{definition}

\subsection{零因子与整环}

\begin{definition}[零因子（Zero Divisor）]
环 $R$ 中若存在 $a\ne 0$，$b\ne 0$ 使得 $ab = 0$，则称 $a$ 为
\keyword{左零因子}、$b$ 为\keyword{右零因子}，统称\keyword{零因子}。
\end{definition}

除环（特别地，域）中没有零因子：若 $ab = 0$ 且 $a\ne 0$，两边左乘
$a^{-1}$ 即得 $b = 0$。

\begin{definition}[整环（Integral Domain）]
无零因子的交换（含幺）环称为\keyword{整环}。
\end{definition}
